% Abstract: 150--200 words (template). Keywords: 4--6, alphabetical, lowercase, not repeating title terms.
\begin{csmdabstract}
Karmada places workloads on member Kubernetes clusters using cluster health, policy and capacity estimates, but it has no notion of the network latency between clients and clusters. This paper designs a cluster selection strategy that adds such latency to Karmada's existing scheduling process without replacing it. An analysis of the Karmada v1.19.1 source code shows four properties that constrain any extension: score plugins are summed without weights, scores influence selection only under a cluster spread constraint, capacity is estimated from requested rather than used resources, and placement is not revisited when cluster conditions change. Based on these properties, the design separates hard eligibility constraints, evaluated by a filter that reuses Karmada's replica estimators, from a soft score that combines post-placement resource headroom with normalised client-to-cluster round-trip time. Latency measurements are attached to cluster objects by an external probe and carry a validity rule: a cluster without a fresh measurement never benefits from the missing value, and the score falls back to resource-only ranking when no candidate has valid data. A worked example and an analysis of weight sensitivity illustrate the decisions of the design. A validation protocol for a future experimental study is proposed.
\end{csmdabstract}

\csmdkeywords{container orchestration, edge-cloud continuum, kubernetes federation, latency, scheduling plugin, telemetry staleness}
